\thispagestyle{empty}

\begin{abstract}
\noindent Излагается доказательство теоремы Нэша о гладком вложении. Упор сделан на простоту и строгость.
\end{abstract}

В римановой геометрии две основные теоремы.
Одна о связности Леви-Чивиты, а другая, доказанная Джоном Нэшем~\cite{nash-1956-ru},
говорит, что \textit{любое риманово многообразие изометрично гладкому подмногообразию евклидова пространства большой размерности}.
То есть римановы многообразия можно определять как подмногообразия евклидовых пространств с индуцированными метриками.

Нэш использовал вариант метода Ньютона.
Однако прямое применение этого метода наталкивается на проблему потери гладкости: поправочные члены требуют контроля большего числа производных, чем дают предыдущие приближения.
Для борьбы с этой трудностью был разработан новый метод, использующий тщательно подобранные сглаживания.
Дальнейшее развитие идеи привело к теореме Нэша--Мозера, нашедшей множество применений.

Позже Маттиас Гюнтер нашёл трюк~\cite{gunter}, устранивший потерю гладкости.
Его доказательство использует сжимающее отображение, ровно как в теореме об обратной функции для банаховых пространств.
Благодаря этому трюку мне удалось включить теорему Нэша в курс дифференциальной геометрии.
По следам этих лекций написана эта заметка, и я надеюсь, что она окажется полезной.

За исключением теоремы об обратной функции, я следую доказательству Нэша, но использую несколько косметических упрощений, предложенных
Дином Янгом~\cite{yang},
Теренсом Тао~\cite{tao} и
Ральфом Ховардом~\cite{howard}.
Подход Гюнтера изложен также в двух учебниках:
Майкла Тейлора~\cite{taylor} и
Цин Ханя и Цзя-Син Хуна~\cite{han-hong}.

Читатель должен уметь работать с гладкими многообразиями (карты, атласы, покрытия, разбиения единицы и тензорные поля).
Другие результаты, включая оценки Шаудера и теорему Уитни о вложении, можно применять как чёрные ящики.

В разделе~\ref{sec1} приводится формулировка задачи о вложении.
В разделе~\ref{sec1+} вводится скручивание Нэша и решается приближенная версия задачи.
Раздел~\ref{sec:Reductions} состоит из нескольких шагов, сводящих задачу вложения к реализации малых возмущений некоторой метрики на торе.
В разделе~\ref{sec:free} вводятся свободные отображения и доказывается лемма Нэша.
В следующем разделе мы формулируем лемму Гюнтера и используем её в доказательстве реализации малых возмущений, а значит, и задачи вложения.
В заключительном разделе доказывается лемма Гюнтера.
Это основная часть статьи, остальное потребовалось, чтобы показать силу этой леммы.

Вот список глобальных обозначений:
\spell{\begin{multicols}{3}}{}
\footnotesize
\begin{list}{}{\leftmargin=.5em}
\item $v\oplus w$, \pageref{oplus}
\item $dw$, $w^*$, $\vec x w$, \pageref{dw}
\item $(d\psi)^2$, \pageref{d2}
\item $\Delta u$, \pageref{Du}
\item $f$, \pageref{f}
\item $L$, $L_{ij}$, $M$, \twopageref{LM}{Lij}
\item $N$, \pageref{Nn}
\item $Q$, $Q_{ij}$, $\tilde Q$, \twopageref{Q}{tildeQ}
\item $\T$, $\T^\perp$, $\T^2$, $\N$, \pageref{TTTN}
\end{list}
\spell{\end{multicols}}{}


\section{Формулировка}\label{sec1}

\paragraph{Индуцированная метрика.}
Пусть $\Omega$ --- гладкое $n$-мерное многообразие.
Его касательное расслоение и касательное пространство в точке $p$ будут обозначаться $\T\Omega$ и $\T_p\Omega$. \index{4t@$\T\Omega$, $\T_p\Omega$}\label{T}
\index{метрический тензор}\emph{Метрический тензор}, скажем $g$, на $\Omega$ --- это гладкое поле симметричных билинейных форм на $\T\Omega$.
Здесь и далее \emph{гладкий} означает~$C^\infty$.

Пусть $\bar\Omega$ --- другое гладкое многообразие, а $w\colon\Omega \to \bar\Omega$ --- гладкое отображение.
Производная $w$ в направлении касательного вектора $\vec x \in \T_p\Omega$ будет обозначаться \label{dw}\index{4d@$w^*$, $dw$, $\vec x w$} $\vec x w = dw(\vec x)$;
формально говоря, производная $w$ в направлении $\vec x$ --- это \index{дифференциал}\emph{дифференциал} $dw\colon\T\Omega \z\to \T\bar\Omega$ отображения $w$, вычисленный на $\vec x$.
Отметим, что $\vec x w \z\in \T_{w(p)}\bar\Omega$.

Предположим, что на $\bar\Omega$ задан метрический тензор $\bar g$.
Тогда метрический тензор $g$ на $\Omega$, заданный равенством
$g(\vec x, \vec y) = \bar g(\vec x w, \vec y w)$,
называется \index{индуцированный метрический тензор}\emph{индуцированным}.
Иначе говоря, $g$ есть поднятие $\bar g$ на $\Omega$ вдоль $w$, что записывается равенством $g = w^*\bar g$.
В этом случае отображение $w\colon(\Omega, g)\to(\bar\Omega, \bar g)$ называется \index{изометрическое отображение}\emph{изометрическим}.

Метрический тензор $g$ называется \index{риманов метрический тензор}\emph{римановым}, если он \textit{положительно определён};
то есть $g(\vec x, \vec x) > 0$ для любого ненулевого касательного вектора $\vec x$.
Гладкое многообразие с заданной римановой метрикой называется \index{риманово многообразие}\emph{римановым многообразием}.

Заметим, что если $g = w^* \bar g$ --- риманова метрика, то $w$ обязано быть регулярным; то есть дифференциал $d_p w$ имеет ранг $n$ в каждой точке $p \in \Omega$.
В частности, $w$ является погружением.
И наоборот, если $w$ регулярно и $\bar g$ риманова, то и $g = w^* \bar g$ риманова.

Пространство $\RR^d$ будет рассмативаться с метрическим тензором, заданным стандартным скалярным произведением:
\[\bar g(\vec x, \vec y) \df \langle \vec x, \vec y \rangle = x_1 \cdot y_1 + \dots + x_d \cdot y_d,\]
где $x_i$ и $y_i$ обозначают координаты векторов $\vec x, \vec y \in \RR^d$.
Заметим, что $\bar g$ является римановой метрикой.

\begin{thm}{Теорема Нэша о вложении}\label{thm:main}
Произвольное $n$-мерное риманово многообразие $(\Omega,g)$ допускает гладкое изометрическое вложение в $\RR^d$, причём $d$ зависит только от~$n$.
\end{thm}

Полное доказательство будет приведено только для компактных многообразий.
Общий случай сводится к компактному, набросок этого доказательства даётся в шаге 4 следующего раздела.

\paragraph{Замечание о размерности \textit{d}.}
Проследив оценки в доказательстве, можно убедиться, что достаточно взять $d = 10 \cdot n^4$; я не пытался оптимизировать эту оценку.

Для компактных поверхностей можно взять $d\z=5$,
и похоже, что это единственный интересный случай с известной оптимальной размерностью.
Существование изометрического вложения в $\RR^5$ доказано в \cite[3.2.4(D)]{gromov-1990-ru}, при этом нет изометрических вложений в $\RR^4$ у $\RP^2$ со стандартной метрикой \cite[приложение 4]{gromov-rokhlin-ru}.

Оценка $d\ge \tbinom{n+1}{2}$ доказана в \cite[приложение 1]{gromov-rokhlin-ru}.
Она означает, что следующая система не может быть переопределённой:
\[
\langle \partial_i w, \partial_j w \rangle = g_{ij}.
\eqlbl{eq:<ww>=g}
\]
Может случиться, что $d=\tbinom{n+1}{2}$ и есть оптимальная оценка при больших $n$ и для локальной версии задачи; то есть \textit{у любой точки гладкого $n$-мерного риманова многообразия есть окрестность $U$, допускающая изометрическое вложение в $\RR^d$}.

Хотя система \ref{eq:<ww>=g} имеет смысл для $C^1$-гладких отображений, доказательство оценки $d\ge \tbinom{n+1}{2}$ существенно использует $C^k$-гладкость для достаточно большого $k$.
Без этого предположения утверждение неверно.
Действительно, по теореме Нэша --- Кёйпера \cite{nash-1954-ru,kuiper-ru} система \ref{eq:<ww>=g} имеет локальное $C^1$-решение при $d = n+1$.

\section{Приближённое решение}\label{sec1+}

\paragraph{Сумма и прямая сумма.}
Пусть $v_1\colon\Omega\z\to\RR^{d_1}$ и $v_2\colon\Omega\z\to\RR^{d_2}$ --- гладкие отображения.
Отображение \label{oplus}\index{$v_1\oplus v_2$}$v_1\oplus v_2\colon p\z\mapsto (v_1(p),v_2(p))$ будет называться \index{прямая сумма}\emph{прямой суммой} $v_1$ и $v_2$;
оно отображает $\Omega$ в $\RR^{d_1+d_2}\z=\RR^{d_1}\oplus\RR^{d_2}$.

Предположим, что метрические тензоры $g$, $g_1$ и $g_2$ индуцированы отображениями $v_1\oplus v_2$, $v_1$ и $v_2$ соответственно.
Тогда
\[g=g_1+g_2.\]

Если $d_1=d_2$, то можно рассмотреть обычную сумму $v_1+v_2$.
Напомним, что \index{носитель}\emph{носитель} отображения --- это замыкание множества, на котором оно принимает ненулевые значения.
Предположим, что $g$, $g_1$ и $g_2$ индуцированы отображениями $v_1+ v_2$, $v_1$ и $v_2$, причём у $v_1$ и $v_2$ непересекающиеся носители.
Тогда снова
\[g=g_1+g_2.\]

Эти два наблюдения будут многократно использоваться далее.

\paragraph{Скручивание Нэша.}
Пусть $\phi$ и $\psi$ --- гладкие функции на гладком многообразии $\Omega$, и $\SSS^1_r$ --- окружность радиуса $r$ в~$\RR^2$.

{

\begin{wrapfigure}[9]{r}{30mm}
\vskip-0mm
\centering
\includegraphics{mppics/pic-15}
\end{wrapfigure}

\mbox{\index{скручивание Нэша}\emph{Скручивание Нэша}} $\Theta\colon\Omega \to \RR^2$ тройки $(r,\phi,\psi)$ определяется как композиция
\[\Omega\xrightarrow{\psi}\RR\xrightarrow{\ell_r}\SSS^1_r\xrightarrow{\times \phi}\RR^2,\]
где $\ell_r\colon\RR\to\SSS^1_r$ --- сохраняющее длину накрытие, например $\ell_r(x)=(r\cdot\cos\tfrac xr,r\cdot\sin\tfrac xr)$,
а $\times \phi$ обозначает умножение на $\phi$; то есть
\[\Theta(x)
\df
\phi(x)\cdot (\ell_r\circ \psi(x)).\]

}

\begin{thm}{Утверждение}\label{clm:twist}
Скручивание Нэша $\Theta$ тройки $(r,\phi,\psi)$ индуцирует метрику
\[g=\phi^2\cdot (d\psi)^2+r^2\cdot(d\phi)^2.\]
Здесь \label{d2}\index{4d@$(d\psi)^2$}$(d\psi)^2$ обозначает метрический тензор, индуцированный функцией
$\psi\colon\Omega \to \RR$; то есть
\[(d\psi)^2(\vec x,\vec y)\df d\psi(\vec x)\cdot d \psi(\vec y)=(\vec x\psi)\cdot (\vec y\psi).\]
\end{thm}

\parit{Вычисление.}
Пусть $\vec x$ --- касательный вектор на $\Omega$. Тогда
\begin{align*}
\vec x\Theta&=\vec x(\phi\cdot \ell_r\circ \psi)=
\\
&=(\vec x\phi)\cdot(\ell_r\circ \psi) + \phi\cdot (\ell_r'\circ \psi)\cdot(\vec x\psi).
\end{align*}

Заметим, что $|\ell_r|=r$, $|\ell'_r|=1$ и $\ell_r\perp \ell'_r$.
Следовательно,
\[g(\vec x,\vec y)=\langle \vec x\Theta,\vec y\Theta\rangle =r^2\cdot (\vec x\phi)\cdot (\vec y\phi)+\phi^2\cdot (\vec x\psi)\cdot (\vec y\psi).
\eqno\qed\]

\paragraph{Приближённое решение.}
Следующее утверждение можно рассматривать как приближённый вариант основной теоремы (\ref{thm:main}).
В частности, из него следует, что множество метрик, индуцированных гладкими отображениями в евклидово пространство, плотно в пространстве всех римановых метрик.

\begin{thm}{Предложение}\label{prop:approx-nash}
Пусть $g$ --- риманова метрика на гладком $n$-мерном многообразии $\Omega$.
Тогда существует метрический тензор $h$ на $\Omega$ и однопараметрическое семейство гладких отображений $w_t\colon\Omega \to \RR^d$ при $t>0$ с индуцированными метриками $g_t=g+t\cdot h$.

Более того, $d$ зависит только от $n$.
\end{thm}

В доказательстве будет построен неотрицательный метрический тензор $h$, но это знание нам не пригодится.

Поскольку $g_t\to g$ при $t \to 0$, возникает соблазн перейти к пределу $w_t$ при $t \to 0$.
Однако, как мы увидим, отображения $w_t$, построенные в доказательстве, сходятся к постоянному отображению, и значит, этот предел задачу не решает.

Пусть $w\colon\Omega \to \RR^d$ --- гладкое отображение.
Предположим, что $w_1,\dots,w_d$ --- координатные функции отображения $w$.
Тогда индуцированный метрический тензор $g$ можно записать в виде
\[g=(dw_1)^2+\dots+(dw_d)^2.
\eqlbl{eq:nash=dw^2}\]
Задача вложения сводится к решению \ref{eq:nash=dw^2}.
Следующая более слабая форма \ref{eq:nash=dw^2} будет играть ключевую роль в доказательстве \ref{prop:approx-nash}.

\begin{thm}{Предложение}\label{prop:phi-f}
Пусть $(\Omega,g)$ --- $n$-мерное риманово многообразие.
Тогда на $\Omega$ существуют гладкие функции
$\phi_1,\z\dots,\phi_d,\psi_1,\z\dots,\psi_d$
такие, что
\[g=\phi_1^2 \cdot (d\psi_1)^2+\dots+\phi_d^2\cdot (d\psi_d)^2.
\eqlbl{eq:g=phidf}\]

Более того, $d$ зависит только от $n$.
\end{thm}


\paragraph{Доказательство.}
Метрический тензор $g$ в локальных координатах $(x_1,\z\dots,x_n)$ можно записать как
\[
g = \sum_{i, j} g_{ij} \cdot dx_i \odot dx_j,
\eqlbl{g=sum}
\]
где $g_{ij} = g_{ji}$ --- гладкие функции от $(x_1,\dots,x_n)$, а $\odot$\label{odot} обозначает симметрическое тензорное произведение.

Поскольку $g$ --- риманова, в любой точке $p \in \Omega$ можно выбрать карту с ортонормированными векторами $\partial_i$ в $p$;
то есть $g_{ii} = 1$ и $g_{ij} = 0$ в точке $p$ для всех $i\ne j$.
Так как $g_{ij}$ --- гладкие функции, для любого $\eps > 0$ найдётся окрестность $U \ni p$ такая, что
\[
g_{ii} \lessgtr 1 \pm \eps
\quad\text{и}\quad
g_{ij}\lessgtr \pm \eps
\eqlbl{alpha=1}\]
при всех $i\ne j$ и в любой точке из $U$.

Заметим, что
\[\pm dx_i\odot dx_j=\tfrac12\cdot (dx_i\pm dx_j)^2-\tfrac12\cdot(dx_i)^2-\tfrac12\cdot(dx_j)^2.\]
Подставим эту формулу в \ref{g=sum} и учтём \ref{alpha=1}.
Если $\eps$ достаточно мало, то $g$ запишется линейной комбинацией $(dx_i)^2$ и $(dx_i\pm dx_j)^2$ с положительными коэффициентами.
Иными словами, в качестве функций $\psi_1,\dots,\psi_d$ можно взять функции $x_i$ и $x_i\pm x_j$ для всех $i<j$; их нужно гладко продолжить на всё многообразие.
Тогда можно найти $\phi_1,\dots,\phi_d$, при которых \ref{eq:g=phidf} выполняется в окрестности точки~$p$.
Применив разбиение единицы, получим глобальное утверждение.

Достаточно взять $n^2$ функций в каждой карте.
Поэтому, если $\Omega$ покрыто $m$ картами, то хватает $d = n^2 \cdot m$.
Более того, та же оценка получается, если карты можно раскрасить в $m$ цветов так, что различные карты одного цвета не пересекаются.
По цветной версии теоремы о размерности Лебега \cite[Theorem 2]{ostrand} можно считать, что $m=n+1$;
то есть хватает $d \z= n^2 \cdot (n+1)$.
\qeds

\paragraph{Доказательство \ref{prop:approx-nash}.}
Выберем функции $\phi_1,\z\dots,\phi_d,\psi_1,\z\dots,\psi_d$ как в \ref{prop:phi-f}.
Обозначим через $\Theta_i$ скручивание Нэша для тройки $(\sqrt{t},\phi_i,\psi_i)$.
По \ref{clm:twist}, семейство
\[w_t=\Theta_1\oplus\dots\oplus \Theta_d\]
является подходящим при
$h=(d\phi_1)^2+\z\dots+(d\phi_d)^2$.
\qedsf

\paragraph{Псевдоевклидово отступление.}
Будем использовать обозначение $\RR^r\ominus\RR^s$\index{3c@$\RR^r\ominus\RR^s$, $v\ominus w$} для псевдоевклидова пространства сигнатуры $(r,s)$;
то есть пространство $\RR^{r+s}$ с метрическим тензором
\[\bar g(\vec x,\vec y)=\langle \vec x,\vec y\rangle=x_1\cdot y_1+\dots+x_r\cdot y_r-x_{r+1}\cdot y_{r+1}-\dots-x_{r+s}\cdot y_{r+s},\]
где $x_i$ и $y_i$ обозначают координаты векторов $\vec x$ и $\vec y$ в $\RR^{r+s}$.

\begin{thm}{Задача}
Покажите, что любой метрический тензор $g$ на гладком многообразии $\Omega$ индуцирован гладким вложением $\Omega\hookrightarrow \RR^r\ominus\RR^s$ при некоторых достаточно больших $r$ и $s$.
\end{thm}

\parit{Решение на основе теоремы Нэша.}
Заметим, что любой метрический тензор $g$ можно записать как разность $g = g_1 - g_2$ двух римановых метрик $g_1$ и $g_2$ на $\Omega$.
По теореме Нэша можно найти два гладких вложения $w_1, w_2 \colon \Omega \hookrightarrow \RR^d$, индуцированные метрики которых равны $g_1$ и $g_2$ соответственно.
Остаётся заметить, что $g$ индуцирован вложением \label{ominus}
\[w_1\ominus w_2\colon p\mapsto(w_1(p),w_2(p))\in \RR^d\ominus\RR^d.\eqno{\qed}\]

Это решение использует теорему вложения, которую мы ещё не доказали.
Следующее решение более прямое; оно использует только скручивание Нэша.

\parit{Прямое решение.}
Предположим, что $g$ --- риманова метрика, и функции $\phi_1,\z\dots,\phi_d,\psi_1,\z\dots,\psi_d$ как в \ref{prop:phi-f}.
Обозначим через $\Theta_i$ скручивание Нэша тройки $(1,\phi_i,\psi_i)$, и пусть
\begin{align*}
h&=(d\phi_1)^2+\dots+(d\phi_d)^2,
\\
\Theta&=\Theta_1\oplus\dots\oplus \Theta_d\colon\Omega\to \RR^{2\cdot d},
\\
\phi&=\phi_1\oplus \dots\oplus \phi_d\colon\Omega\to \RR^{d}.
\end{align*}
Согласно \ref{clm:twist}, $\Theta$ индуцирует метрику $g+h$.
Заметим, что $\phi$ индуцирует $h$.
Значит, $w=\Theta\ominus \phi\colon\Omega\to \RR^{2\cdot d}\ominus\RR^d$ индуцирует $g$.

В общем случае представим $g$ как разность двух римановых метрик: $g = g_1 - g_2$.
Из вышесказанного, получаем гладкие отображения $w_1$ и $w_2$ в псевдоевклидовы пространства, индуцирующие $g_1$ и $g_2$ соответственно.
Отсюда следует, что $w_1 \ominus w_2$ индуцирует $g$.

Наконец, по теореме Уитни можно выбрать вложение $w_3\colon\Omega\hookrightarrow \RR^{2\cdot n+1}$.
Остаётся заметить, что $(w_1 \oplus w_3)\ominus (w_2\oplus w_3)$ --- искомое гладкое изометрическое вложение.
\qeds

\section{Упрощение задачи}\label{sec:Reductions}

\paragraph{Реализация малых возмущений.}
Пусть $\TT^n=\RR^n/\ZZ^n$ --- $n$-мерный тор.
Следующая теорема будет доказана в последующих разделах;
заметим, что это частный случай основной теоремы.

\begin{thm}{Теорема}\label{thm:restricted}
Существует риманова метрика $g_0$ на $\TT^n$ такая, что
любая метрика $g$, достаточно близкая к $g_0$ в $C^\infty$-топологии, индуцируется гладким отображением $\TT^n\to \RR^d$ при некотором $d$, зависящем только от $n$.
\end{thm}

Мы сведём основную теорему (\ref{thm:main}) к этому частному случаю за четыре шага.

\paragraph{Шаг 1.}
Сейчас мы выведем из \ref{thm:restricted} существование изометрического погружения для любой метрики на торе; а именно,
\begin{clm}{}
для любой римановой метрики $g$ на $\TT^n$ существует изометрическое погружение $(\TT^n, g) \looparrowright \RR^d$, причём $d$ зависит только от~$n$.
\end{clm}

Действительно, пусть $g_0$ --- метрика из \ref{thm:restricted}.
Умножив её на достаточно малую положительную константу, можно считать, что $g_0<g$; то есть $g_0(\vec x,\vec x)<g(\vec x,\vec x)$ для любого ненулевого касательного вектора $\vec x$.
В частности, разность $g_1=g-g_0$ является римановой метрикой.

По \ref{prop:approx-nash}, существует метрический тензор $h$ такой, что для любого $t>0$ найдётся гладкое отображение $w_1\colon\TT^n\to \RR^{d_1}$, индуцирующее $g_1+t\cdot h$.

По \ref{thm:restricted}, для любого достаточно малого $t>0$ найдётся гладкое отображение $w_0\colon\TT^n\to \RR^{d_0}$, индуцирующее метрику $g_0-t \cdot h$.
Следовательно, $w_0\oplus w_1$ индуцирует
\[g=g_0-t\cdot h+g_1+t\cdot h.\]

\paragraph{Шаг 2.}
Теперь усилим утверждение до вложения; а именно,
\begin{clm}{}
для любой римановой метрики $g$ на $\TT^n$ существует изометрическое вложение $(\TT^n,g)\hookrightarrow\RR^d$, причём $d$ зависит только от~$n$.
\end{clm}

Выберем гладкое регулярное вложение $w_1\colon\TT^n\hookrightarrow \RR^{n+1}$.
Пусть $g_1$ --- метрика, индуцированная $w_1$.
Так как $w_1$ регулярно, $g_1$ должна быть римановой.

Можно считать, что $g>g_1$; иными словами, метрика $g_2=g-g_1$ является римановой.
Если это не так, то заменим $w_1$ на $\varepsilon \cdot w_1$ при достаточно малом $\varepsilon > 0$.

Шаг 1 даёт изометрическое погружение $w_2\colon(\TT^n,g_2)\z\looparrowright \RR^d$.
Так как $g=g_1+g_2$,
\[w=w_1\oplus w_2\colon (\TT^n,g)\to \RR^{d+ n+1}\]
является изометрическим вложением.
Действительно, $w$ изометрично, и оно является вложением, поскольку $w_1$ --- вложение.

\paragraph{Шаг 3.}
Теперь докажем вложимость всех компактных многообразий; а именно,
\begin{clm}{}
любое компактное $n$-мерное риманово многообразие $(\Omega,g)$ допускает гладкое изометрическое вложение в $\RR^d$, причём $d$ зависит только от~$n$.
\end{clm}

Теорема Уитни даёт гладкое вложение $\iota\colon\Omega\hookrightarrow\TT^{2\cdot n+1}$.

Риманову метрику на $\Omega$ можно продолжить до римановой метрики на торе.
Например, для этого можно выбрать метрику на трубчатой окрестности, которая ограничивается на $\Omega$ в данную метрику, а затем, воспользовавшись разбиением единицы, склеить её со стандартной метрикой на торе.
Остаётся взять гладкое изометрическое вложение $w\colon (\TT^{2\cdot n+1}, \bar g) \hookrightarrow \RR^d$ из шага 2 и заметить, что композиция $w \circ \iota\colon(\Omega, g) \hookrightarrow \RR^d$ есть изометрическое вложение.

\paragraph{Шаг 4.}
Теперь наметим, как свести общий случай к компактному.
Полное доказательство приводится в \cite[Часть D]{nash-1956-ru}.

Рассуждая как на втором шаге, можно свести задачу к нахождению изометрического погружения $(\Omega,g)$.

Далее, строится последовательность отображений $r_{\mathfrak{i}}\colon\Omega \z\to \SSS^n$ и последовательность римановых метрик $h_{\mathfrak{i}}$ на $\SSS^n$ таких, что
\[g=\sum_{\mathfrak{i}} r_{\mathfrak{i}}^*h_{\mathfrak{i}}\]
и $r_{\mathfrak{i}}$ постоянно, скажем равно $c_{\mathfrak{i}}$, вне компактного подмножества $K_{\mathfrak{i}}\subset \Omega$.
Эта часть делается с помощью разбиения единицы.

Кроме того, нужно устроить конструкцию так, чтобы множество индексов $\mathfrak{I}$ можно было раскрасить в $n+1$ цвет так, что $K_{\mathfrak{i}}\cap K_{\mathfrak{j}}=\emptyset$ для различных индексов $\mathfrak{i}$ и $\mathfrak{j}$ одного цвета.
Точнее, $\mathfrak{I}$ можно разбить на $n+1$ подмножество $\mathfrak{I}_0, \dots, \mathfrak{I}_n$ так, что для любого $k$ и любых различных $\mathfrak{i}, \mathfrak{j} \z\in \mathfrak{I}_k$ имеем
$K_{\mathfrak{i}} \cap K_{\mathfrak{j}} = \emptyset$.
Это очень похоже на цветную версию теоремы о размерности Лебега.

Применив компактный случай теоремы к $(\SSS^n,h_{\mathfrak{i}})$, получим изометрическое вложение $w_{\mathfrak{i}}\colon(\SSS^n,h_{\mathfrak{i}})\hookrightarrow \RR^d$.
Сдвинув $w_{\mathfrak{i}}$, можно считать, что $w_{\mathfrak{i}}(c_{\mathfrak{i}})=0$, а значит $w_{\mathfrak{i}}\circ r_{\mathfrak{i}}$ имеет носитель в $K_{\mathfrak{i}}$.
Теперь заметим, что отображение
\[\biggl(\,\sum_{\mathfrak{i}\in \mathfrak{I}_0}w_{\mathfrak{i}}\circ r_{\mathfrak{i}}\biggr)\oplus\dots\oplus
\biggl(\,\sum_{\mathfrak{i}\in \mathfrak{I}_n}w_{\mathfrak{i}}\circ r_{\mathfrak{i}}\biggr)\]
даёт требуемое изометрическое погружение $(\Omega,g)\looparrowright \RR^{(n+1)\cdot d}$.

\section{Свободные отображения}\label{sec:free}

С этого момента рассматривается только случай тора $\TT^n\z=\RR^n/\ZZ^n$, он снабжён циклическими координатами и глобально определёнными частными производными $\partial_1,\dots,\partial_n$.
Кроме того, метрический тензор на $\TT^n$ задаётся гладким отображением $g\colon\TT^n\z\to\RR^N$ с компонентами $g_{ij}$ при $1\le i\le j\le n$; здесь и далее \label{Nn}
\[N\df\tfrac{n\cdot(n+1)}2.\]

\paragraph{Свободные отображения.}
Напомним, что гладкое отображение $f\colon\TT^n \to \RR^d$ \index{регулярное отображение}\emph{регулярно}, если его дифференциал $df$ имеет ранг $n$ в каждой точке.
Иными словами, в любой точке $p\in \TT^n$ частные производные первого порядка $\partial_1 f,\dots,\partial_n f$ линейно независимы.

Отображение $f\colon\TT^n \to \RR^d$ называется \index{свободное отображение}\emph{свободным}, если аналогичное свойство выполняется для частных производных первого и второго порядка:
то есть в любой точке $p\in \TT^n$ все $n + N$ частных производных
$\partial_i f$,
$\partial_{ij}f$ при $i\le j$
линейно независимы.

Пусть $\psi_\ell\colon\TT^n\to\SSS^1=\RR/\ZZ$ --- отображение, индуцированное линейной функцией $\ell\colon\RR^n\to\RR$.
Будем рассматривать $\SSS^1$ как подмножество в $\RR^2$.
Пусть $f\colon\TT^n\z\to \RR^{n\cdot(n+1)}$ --- прямая сумма отображений
$\psi_{x_i}$ и $\psi_{x_i+x_j}$, где $x_i$ --- координатные функции на $\RR^n$ и $i<j$.
Нетрудно проверить, что $f$ свободно; тем самым доказано следующее утверждение.

\begin{thm}{Утверждение}\label{ex:free-composition}
Существует гладкое свободное вложение $f\colon\TT^n\hookrightarrow \RR^{n\cdot(n+1)}$.
\end{thm}

С этого момента свободное отображение $f\colon\TT^n\hookrightarrow \RR^d$\label{f} зафиксировано, и
тор $\TT^n$ рассматривается как подмногообразие в $\RR^d$.
Касательное пространство $\T_p=\T_p\TT^n$ --- это $n$-мерное пространство, натянутое на частные производные $\partial_1f,\z\dots, \partial_nf$ в~$p$, а $\T_p^\perp$ --- его ортогональное дополнение в $\RR^d$.
Соприкасающееся пространство $\T^2_p$ определяется как линейная оболочка частных производных первого и второго порядка $\partial_i f$ и $\partial_{ij} f$ в точке $p$, и $\N_p \z\df \T^2_p \cap \T_p^\perp$.

Так как $f$ свободно, размерности пространств $\T_p$, $\T_p^\perp$, $\T^2_p$ и $\N_p$ равны $n$, $d - n$, $n + N$ и $N$ соответственно.
Мы будем писать $\T$, $\T^\perp$, $\T^2$ и $\N$ для соответствующих векторных расслоений над $\TT^n$.
\index{4t@$\T$, $\T^\perp$, $\T^2$, $\N$}\label{TTTN}

\paragraph{Нормы Гёльдера.}
Выберем раз и навсегда некоторое $\alpha\in(0,1)$, например $\alpha=\tfrac12$.
Через $|x-y|_{\TT^n}$ будем обозначать стандартное расстояние между точками $x,y\z\in\TT^n$.
Напомним определение норм Гёльдера на открытом подмножестве $W\subset \TT^n$: \label{C(T)}
\begin{align*}
\|u\|_{C^{0,\alpha}(W)}
&\df\sup_{x}\{\,|u(x)|\,\}
+\sup_{x\ne y}\left\{\,\frac{|u(x)-u(y)|}{|x-y|_{\TT^n}^\alpha}\,\right\},
\\
\|u\|_{C^{k,\alpha}(W)}&\df\max_{|\mathcal{I}|\le k}\left\{\,\|\partial^{\mathcal{I}}u\|_{C^{0,\alpha}(W)}\,\right\},
\end{align*}
где $x,y\in W$, а $\mathcal{I} = (i_1,\dots,i_n)$ обозначает мультииндекс; таким образом, $\partial^{\mathcal{I}}$ --- это частная производная $\partial_1^{i_1}\dots\partial_n^{i_n}$ порядка $|\mathcal{I}| \z\df i_1 +\dots+i_n$.
Пространство функций $W\to \RR$ с конечной ${C^{k,\alpha}(W)}$-нормой называется $(k,\alpha)$-пространством Гёльдера на $W$ и обозначается $C^{k,\alpha}(W)$; \index{$C^{k,\alpha}$, $\Arrowvert u\Arrowvert_{C^{k,\alpha}}$} в этих обозначениях можно опускать $W$.

Пусть $C^{k,\alpha}(\TT^n,\RR^d)$ --- пространство отображений $u\colon\TT^n\z\to \RR^d$, с координатными функциями $u_1,\dots, u_d$ из $C^{k,\alpha}(\TT^n)$;
положим
\[\|u\|_{C^{k,\alpha}}\df\max \{\,\|u_1\|_{C^{k,\alpha}},\dots,\|u_d\|_{C^{k,\alpha}} \,\}.\]

\paragraph{Q-форма.}
Нам потребуется немного причудливое обозначение для индуцированной метрики.
Определим симметрическую билинейную форму $Q$\index{1Q@$Q(v,w)$, $Q_{ij}(v,w)$}\label{Q}, которая по гладким отображениям $v, w\colon\TT^n \z\to \RR^d$ выдаёт следующий метрический тензор на~$\TT^n$:
\[
g(\vec x, \vec y)
\df
\tfrac{1}{2} \cdot [\langle \vec x v, \vec y w \rangle + \langle \vec x w, \vec y v \rangle].
\]
Напомним, что метрический тензор на $\TT^n$ записывается в циклических координатах как гладкое отображение $\TT^n\to\RR^N$.
Компоненты метрического тензора записываются как
\[
g_{ij} = Q_{ij}(v, w) = \tfrac{1}{2} \cdot [\langle \partial_i v, \partial_j w \rangle+\langle \partial_i w, \partial_j v \rangle].
\]

Заметим, что метрический тензор $g$ индуцирован гладким отображением $w\colon\TT^n \z\to \RR^d$
тогда и только тогда, когда $g = Q(w, w)$.


\paragraph{Линеаризация.}
Напомним, что $f\colon\TT^n \hookrightarrow\RR^d$ --- фиксированное гладкое свободное вложение
(оно существует по \ref{ex:free-composition}).
Определим линейный оператор
\[L\colon C^\infty(\TT^n,\RR^d)\to C^\infty(\TT^n,\RR^N)\qquad
\text{как}
\qquad L(w)\df 2\cdot Q(f,w).\]

Скоро мы увидим, что реализация малых возмущений метрики следует из разрешимости квадратного уравнения
\[L(w)+Q(w,w)\z=h\]
при достаточно малых $h$.
Следующая лемма даёт решение соответствующего линейного уравнения $L(w)=h$.

\begin{thm}{Лемма Нэша}\label{lem:nash-guenter:nash}
$L$ допускает правый обратный оператор \index{1Q@$L(w)$, $M(h)$}\label{LM}
\[M\colon C^\infty(\TT^n,\RR^N)\to C^\infty(\TT^n,\RR^d)\]
такой, что
\[\|Mh\|_{C^{2,\alpha}}\le a\cdot \|h\|_{C^{2,\alpha}}\]
для некоторой константы $a$ и любого $h\in C^\infty(\TT^n,\RR^N)$.
Более того, $Mh$ является $\N$-полем;
то есть $Mh(p)\in \N_p\z=\T_p^\perp\cap \T_p^2$ для всех $p\in \TT^n$.

\end{thm}

\paragraph{Доказательство.}
Выберем метрический тензор $h$ на $\TT^n$.
Так как все $\partial_i f$,
$\partial_{ij}f$
линейно независимы,
уравнения
$\langle\partial_{ij}f,w\rangle=-\tfrac12\cdot h_{ij}$
однозначно задают $\N$-поле $w$.
Поле $w$ можно также определить как линейную комбинацию $\partial_i f$,
$\partial_{ij}f$ такую, что
\[
\begin{aligned}
\langle\partial_if,w\rangle&=0,
&
-2\cdot\langle\partial_{ij}f,w\rangle&=h_{ij}
\end{aligned}
\eqlbl{eq:w(h)}
\]
в каждой точке.

Заметим, что \label{Lij}
\begin{align*}
L_{ij}(w)&=2\cdot Q_{ij}(f,w)=
\\
&=\langle\partial_if,\partial_jw\rangle+\langle\partial_jf,\partial_iw\rangle=
\\
&=\partial_j\langle\partial_if,w\rangle+\partial_i\langle\partial_jf,w\rangle -2\cdot \langle\partial_{ij}f,w\rangle.
\end{align*}
Согласно \ref{eq:w(h)}, $M\colon h \mapsto w$ --- правый обратный оператор к~$L$.
\qeds

\section{Реализация малых возмущений}\label{sec:perturbation}

\begin{thm}{Лемма Гюнтера}\label{lem:nash-guenter:guenter}
Существует симметрическая билинейная форма \index{1Q@$\tilde Q(v,w)$}\label{tildeQ}
\[\tilde Q\colon C^\infty(\TT^n,\RR^d)\times C^\infty(\TT^n,\RR^d) \to C^\infty(\TT^n,\RR^d)\]
такая, что $L\tilde Q=Q$ и для некоторой последовательности констант $b_2,b_3,\dots$ выполняются неравенства
\begin{align*}
\|\tilde Q(w,w)\|_{C^{2,\alpha}}&\le b_2\cdot \|w\|_{C^{2,\alpha}}^2,
\\
\|\tilde Q(w,w)\|_{C^{k,\alpha}}
&\le b_2\cdot \|w\|_{C^{2,\alpha}}\cdot \|w\|_{C^{k,\alpha}}
+
b_k\cdot \|w\|_{C^{k-1,\alpha}}^2
\end{align*}
при любом $w\in C^\infty(\TT^n,\RR^d)$ и $k\ge 3$.
\end{thm}

Доказательство приводится в следующем разделе.

\paragraph{Реализация малых возмущений.}
Сейчас мы завершим доказательство теоремы Нэша, по модулю леммы Гюнтера.
Напомним, что $f\colon\TT^n \hookrightarrow\RR^d$ --- фиксированное гладкое свободное вложение и $L(w)\z\df 2\cdot Q(f,w)$.

\begin{thm}{Лемма}\label{lem:perturb}
Существует $r>0$ такое, что уравнение
\[L(w)+Q(w,w)=h\eqlbl{eq:L+Q=h}\]
имеет гладкое решение $w\colon\TT^n\to\RR^d$ для любого гладкого метрического тензора~$h$ (не обязательно риманова), удовлетворяющего $\|h\|_{C^{2,\alpha}}\z<r$.
Более того, можно считать, что
\[\|w\|_{C^{2,\alpha}}\le c\cdot \|h\|_{C^{2,\alpha}}\]
для некоторой константы $c$.
\end{thm}

В разделе \ref{sec:Reductions} мы показали, что реализация малых возмущений (\ref{thm:restricted}) влечёт основную теорему (\ref{thm:main}).
Теперь покажем, что лемма \ref{lem:perturb} влечёт реализацию малых возмущений (\ref{thm:restricted}), а значит, и основную теорему (\ref{thm:main}).

\paragraph{Доказательство \ref{thm:restricted} по модулю \ref{lem:perturb}.}
Напомним, что $f\colon\TT^n\z\to\RR^d$ --- фиксированное свободное отображение.
Пусть $g_0$ --- метрический тензор, индуцированный $f$, то есть $g_0=Q(f,f)$.
Далее, положим $h=g-g_0$,
тогда $h$ --- малый метрический тензор; в частности, можно считать, что $\|h\|_{C^{2,\alpha}}\z<r$.

Достаточно найти гладкое отображение $w\colon\TT^n\to\RR^d$ такое, что
\[
Q(f+w, f+w) = g_0 + h,
\]
а это уравнение эквивалентно \ref{eq:L+Q=h}.
Действительно,
\begin{align*}
Q(f+w, f+w)
&= Q(f, f) + 2 \cdot Q(f, w) + Q(w, w)
\\
&= g_0 + L(w) + Q(w, w).
\end{align*}
Следовательно, \ref{thm:restricted} доказана. \qeds

Доказательство тоже, что в классической теореме об обратной функции в банаховых пространствах; см., например, \cite{cartan-ru}.

\paragraph{Доказательство \ref{lem:perturb} по модулю \ref{lem:nash-guenter:guenter}.}
Напомним, что $M$, $\tilde Q$, $a$ и $b_2$, $b_3,\dots$ определены в \ref{lem:nash-guenter:nash} и \ref{lem:nash-guenter:guenter}.
Положим $R=\tfrac1{10\cdot b_2}$ и $r=\tfrac R{10\cdot a}$.
Предположим, что $\|h\|_{C^{2,\alpha}}\z\le r$.
Согласно \ref{lem:nash-guenter:guenter},
\[\Phi\colon w\mapsto Mh-\tilde Q(w,w)\]
--- сжимающее отображение шара $\cBall[0,R]\subset {C^{2,\alpha}(\TT^n,\RR^d)}$ в себя.

Рассмотрим последовательность отображений $w_0,w_1,\dots\colon\TT^n\z\to\RR^d$, заданную условиями $w_0=0$ и $w_{n+1}=\Phi(w_n)$ для всех $n$.
Заметим, что каждое $w_n$ гладко.
Так как $\Phi$ сжимающее, последовательность $w_n$ сходится в $C^{2,\alpha}(\TT^n,\RR^d)$;
пусть $w$ --- её предел.
Тогда
\[\|w\|_{C^{2,\alpha}}\le R
\qquad\text{и}\qquad
w= Mh-\tilde Q(w,w).
\eqlbl{eq:w=Mh-Q(w,w)}\]

Так как $LM=\id$ и $L\tilde Q=Q$, применив $L$ к обеим частям равенства \ref{eq:w=Mh-Q(w,w)}, получим, что $w$ удовлетворяет \ref{eq:L+Q=h}.

Воспользовавшись следующим утверждением, покажем, что $w$ гладкое:
\begin{clm}{}\label{clm:wn-bounded}
Нормы $\|w_n\|_{C^{k,\alpha}}$ ограничены для каждого целого $k\ge 2$.
\end{clm}
Действительно, вложение
$C^{k,\alpha}(\TT^n,\RR^d)\hookrightarrow C^k(\TT^n,\RR^d)$ компактно.
Поэтому найдётся подпоследовательность $w_n$, сходящаяся в $C^k(\TT^n,\RR^d)$,
и её предел обязан совпасть с $w$.
В частности $w$ является $C^k$-гладким для каждого $k$, а значит гладким.

Утверждение \ref{clm:wn-bounded} будет доказано индукцией по $k$.
База $k=2$ уже есть.

По предположению индукции, нормы $\|w_n\|_{C^{k-1,\alpha}}$ ограничены.
В частности, существует константа $c_k$ такая, что
\[\|Mh\|_{C^{k,\alpha}}+b_k\cdot \|w_n\|_{C^{k-1,\alpha}}^2\le c_k\]
для всех $n$; здесь $b_k$ --- константа из \ref{lem:nash-guenter:guenter}.

Так как $\|w_n\|_{C^{2,\alpha}}\le R=\tfrac1{10\cdot b_2}$, из второй оценки в \ref{lem:nash-guenter:guenter} получаем
\begin{align*}
\|w_{n+1}\|_{C^{k,\alpha}}
&\le
\|Mh\|_{C^{k,\alpha}}
+
b_2\cdot \|w_n\|_{C^{2,\alpha}}\cdot \|w_n\|_{C^{k,\alpha}}
+
b_k\cdot \|w_n\|_{C^{k-1,\alpha}}^2\le
\\
&\le c_k +\tfrac1{10}\cdot\|w_n\|_{C^{k,\alpha}}.
\end{align*}
В частности, $\|w_n\|_{C^{k,\alpha}}\le 2\cdot c_k
\ \Rightarrow\ \|w_{n+1}\|_{C^{k,\alpha}}\le 2\cdot c_k$
для всех $n$.
Так как $w_0=0$, получаем \ref{clm:wn-bounded}.
\qeds

\section{Доказательство леммы Гюнтера}\label{sec:proof(gunther)}

Напомним, что $L$ и $M$ определены на странице \pageref{LM}.
Форма $\tilde Q$ в лемме Гюнтера определяется равенством
\[\tilde Q(w,w)=M(Q(w,w)+L\vec x)-\vec x,\]
где $\vec x$ --- касательное векторное поле на $\TT^n$, квадратично зависящее от $w$;
см. замечание сразу после доказательства ниже.
Так как $M$ является правым обратным к $L$, имеем $L\tilde Q=Q$.
Трюк состоит в том, чтобы выбрать $\vec x$ так, чтобы выполнились неравенства в лемме; см. \ref{eq:Delta}, \ref{eq:AB=<C} и \ref{eq:tildeQ} ниже.

В доказательстве \ref{lem:perturb} мы нашли поправку $w$ к свободному вложению $f$, удовлетворяющую равенству
\[w= Mh-\tilde Q(w,w)=M(h-Q(w,w)-L\vec x)+\vec x.\]
По лемме Нэша, $\vec x$ является касательной частью $w$.
Это касательное слагаемое --- главное отличие от оригинальной схемы Нэша.
Заметим, что возмущение вложения в касательном направлении почти не меняет риманову метрику --- в первом порядке оно лишь меняет параметризацию.
По этой причине введение касательного слагаемого может выглядеть странно, но оно работает.
Та же идея срабатывает в трюке Детурка для потока Риччи \cite[2.2]{brendle} и в уточнении теоремы Нэша о $C^1$-вложении Николаса Кёйпера~\cite{kuiper-ru}.

\paragraph{Сглаживающий оператор.}
Напомним, что $\alpha\in(0,1)$ фиксировано, например $\alpha=\tfrac12$.

Лапласиан гладкой функции $u\colon \TT^n\to\RR$ определяется формулой \index{4d@$\Delta u$}\label{Du}
\[\Delta u=\sum_s\partial_{ss}u,\]
где $\partial_1,\dots,\partial_n$ --- стандартные частные производные на $\TT^n$.

\begin{thm}{Предложение}\label{prop:Delta-1}
У оператора $(\Delta-1)\colon u\mapsto \Delta u-u$ на $C^\infty(\TT^n)$ существует обратный $(\Delta-1)^{-1}$.
Более того, существует константа $c>0$ такая, что при любом целом $k\ge2$ имеем
\begin{align*}
\|(\Delta-1)u\|_{C^{k-2,\alpha}(\TT^n)}&\le c\cdot \|u\|_{C^{k,\alpha}(\TT^n)}
\\
\|(\Delta-1)^{-1}v\|_{C^{k,\alpha}(\TT^n)}&\le c\cdot \|v\|_{C^{k-2,\alpha}(\TT^n)}
\end{align*}
для любых функций $u,v\in C^\infty(\TT^n)$.
\end{thm}

\parit{Набросок доказательства.}
Существование $(\Delta-1)^{-1}$ вытекает из анализа Фурье.
Иначе говоря, уравнение $\Delta u - u = v$ имеет единственное гладкое решение $u$ при данном $v\in C^\infty(\TT^n)$.

Остаётся доказать следующее неравенство с некоторой константой $c$:
\[\tfrac1c\cdot\|v\|_{C^{k-2,\alpha}(\TT^n)}
\le
\|u\|_{C^{k,\alpha}(\TT^n)}
\le
c\cdot\|v\|_{C^{k-2,\alpha}(\TT^n)};
\eqlbl{eq:biLip}\]
первое неравенство тривиально.

По принципу максимума имеем
\[\|u\|_{L^\infty(\TT^n)}
\le
\|v\|_{L^\infty(\TT^n)},
\eqlbl{eq:maximum-principle}\]
где $\|\ \|_{L^\infty(W)}$ обозначает $\sup$-норму на $W\subset \TT^n$.

Выберем $\tfrac1{10}$-шар $B$ в $\TT^n$ со стандартной метрикой;
пусть $2\cdot B$ --- концентрический шар вдвое большего радиуса.
Так как $2\cdot B$ изометричен евклидову шару, к сужению $u$ на $2\cdot B$ применима оценка Шаудера; см., например, \cite[лемма 6.1(a)]{gilbarg-trudinger-ru}.
Применив \ref{eq:maximum-principle}, получаем
\begin{align*}
\|u\|_{C^{2,\alpha}(B)}
&\le c\cdot(\|v\|_{C^{0,\alpha}(2\cdot B)} + \|u\|_{L^\infty(2\cdot B)})\le
2\cdot c\cdot\|v\|_{C^{0,\alpha}(\TT^n)}.
\end{align*}
Усреднив это неравенство по $\tfrac1{10}$-шарам в $\TT^n$, получим второе неравенство в \ref{eq:biLip} при $k=2$.

Поскольку $\partial_i\Delta=\Delta\partial_i$ для каждого $i$,
применив уже доказанный случай к частным производным производным $\partial^{\mathcal{I}}u$ и $\partial^{\mathcal{I}}v$, поручаем \ref{eq:biLip} при всех $k$.
\qeds

Заметим, что $(\Delta-1)^{-1}$ --- сглаживающий оператор.
Более того, $(\Delta-1)^{-1}$ коммутирует с частными производными; то есть
\[(\Delta-1)^{-1}\partial_i
=
\partial_i(\Delta-1)^{-1}.\]
Эти два свойства будут играть ключевую роль в следующем доказательстве.

Кроме того, нам понадобится следующее утверждение, вытекающее из правила произведения для производных.

\begin{thm}{Утверждение}\label{ex:hoelder-norm}
Для данной функции $v\in C^\infty(\TT^n)$ существует последовательность констант $a_0,a_1,\dots$ такая, что
\[\|u\cdot v\|_{C^{k,\alpha}}\le a_0\cdot\|u\|_{C^{k,\alpha}}+a_k\cdot\|u\|_{C^{k-1,\alpha}}\]
для любого целого $k\ge 1$ и любого $u\in C^\infty(\TT^n)$.
\end{thm}

\paragraph{Доказательство \ref{lem:nash-guenter:guenter}.}
Напомним, что
$Q_{ij}(w,w)=\langle\partial_i w,\partial_jw\rangle$.
Покажем, что
\[(\Delta-1) Q_{ij}(w,w)=A_{ij}(w)
+\partial_iA_j(w)
+\partial_j A_i(w),\eqlbl{eq:Delta}\]
где $A_{ij}$, $A_{i}$ и $A_{j}$ --- линейные комбинации (с постоянными коэффициентами) следующих квадратичных форм:
\[
\begin{aligned}
w&\mapsto \langle\partial_aw,\partial_bw\rangle,
&
w&\mapsto \langle\partial_{ab}w,\partial_cw\rangle,
&
w&\mapsto \langle\partial_{ab}w,\partial_{cd}w\rangle.
\end{aligned}
\eqlbl{eq:forms}
\]
Здесь $a,b,c,d$ пробегают все индексы.

Действительно,
\begin{align*}
(\Delta-1) Q_{ij}(w,w)&=
\langle\partial_{i}\Delta w,\partial_jw\rangle
+\langle\partial_{i} w,\partial_{j}\Delta w\rangle+
\\
&\qquad\qquad+2\cdot \sum_s\langle\partial_{is} w,\partial_{js}w\rangle
-\langle\partial_{i} w,\partial_{j}w\rangle
=
\\
&=
2\cdot \sum_s
\langle\partial_{is} w,\partial_{js}w\rangle
-\langle\partial_{i} w,\partial_{j}w\rangle
-2\cdot \langle\Delta w,\partial_{ij}w\rangle
+
\\
&\qquad\qquad+\partial_i\langle\Delta w,\partial_{j}w\rangle
+\partial_j\langle\Delta w,\partial_{i} w\rangle,
\end{align*}
что и доказывает \ref{eq:Delta}.

Теперь выберем $A=A_i$ или $A=A_{ij}$ при некоторых $i$ и $j$.
Заметим, что существует последовательность констант $c_2,c_3,\dots$ такая, что
\[
\begin{aligned}
\|A(w)\|_{C^{0,\alpha}}&\le c_2\cdot \|w\|^2_{C^{2,\alpha}},
\\
\|A(w)\|_{C^{k-2,\alpha}}&\le c_2\cdot \|w\|_{C^{k,\alpha}}\cdot \|w\|_{C^{2,\alpha}}+c_{k}\cdot\|w\|^2_{C^{k-1,\alpha}}.
\end{aligned}
\eqlbl{eq:A=<C}\]

Эти неравенства получаются проверкой каждой формы из \ref{eq:forms}.
Разберём случай $w\mapsto \langle\partial_{ab}w,\partial_{cd}w\rangle$.
Первое неравенство в \ref{eq:A=<C} тривиально.
Второе следует из того, что для любой частной производной $\partial^{\mathcal{I}}$ порядка $k-2$ имеем
\begin{align*}
\partial^{\mathcal{I}}\langle\partial_{ab}w,\partial_{cd}w\rangle&=\langle\partial_{ab}\partial^{\mathcal{I}}w,\partial_{cd}w\rangle+\langle\partial_{ab}w,\partial_{cd}\partial^{\mathcal{I}}w\rangle+\\
&\qquad\qquad + \text{дополнительные слагаемые},
\end{align*}
где каждое дополнительное слагаемое имеет вид $\langle\partial^{\mathcal{K}}w,\partial^{\mathcal{L}}w\rangle$, а частные производные
$\partial^{\mathcal{K}}$ и $\partial^{\mathcal{L}}$ имеют порядок не больше $k-1$.

Согласно \ref{prop:Delta-1}, можно выбрать новую последовательность констант $c_2, c_3, \dots$ такую, что
\[
\begin{aligned}
\|(\Delta-1)^{-1}A(w)\|_{C^{2,\alpha}}
&\le c_2\cdot \|w\|_{C^{2,\alpha}}^2,
\\
\|(\Delta-1)^{-1}A(w)\|_{C^{k,\alpha}}
&\le c_2\cdot \|w\|_{C^{k,\alpha}}\cdot \|w\|_{C^{2,\alpha}} +c_{k}\cdot\|w\|_{C^{k-1,\alpha}}^2.
\end{aligned}
\eqlbl{eq:AB=<C}\]

Напомним, что
\[
L_{ij}(v)=-2\cdot\langle\partial_{ij} f,v\rangle
+\partial_{i}\langle\partial_{j} f,v\rangle
+\partial_{j}\langle\partial_{i} f,v\rangle;
\eqlbl{eq:L-def2}
\]
см. доказательство леммы Нэша.
Определим
$\tilde Q(w,w)$ как линейную комбинацию $\partial_if$ и $\partial_{ij}f$ для всех $i\le j$
такую, что
\[
\begin{aligned}
-2\cdot\langle \partial_{ij}f,\tilde Q(w,w)\rangle
&=(\Delta-1)^{-1} A_{ij}(w),
\\
\langle \partial_if,\tilde Q(w,w)\rangle
&=(\Delta-1)^{-1} A_i(w)
\end{aligned}
\eqlbl{eq:tildeQ}
\]
при всех $i\le j$.
Так как $f$ свободно, $\tilde Q(w,w)$ корректно определено и продолжается до симметрической билинейной формы.
Объединяя \ref{eq:L-def2}, \ref{eq:tildeQ} и \ref{eq:Delta}, получаем
\[
L\tilde Q=Q.
\]

Пусть $\{e_i,e_{ij}\}_{i\le j}$ --- двойственный репер к $\{\partial_if,\partial_{ij}f\}_{i\le j}$ в соприкасающемся расслоении $\T^2$.
Заметим, что
\[\tilde Q(w,w)=\sum_i(\Delta-1)^{-1} A_i(w)\cdot e_i-\tfrac12\cdot\sum_{i\le j}(\Delta-1)^{-1}A_{ij}(w)\cdot e_{ij}.\eqlbl{eq:tildeQend}\]
Так как отображения $e_i, e_{ij}\colon \TT^n \to \RR^d$ гладкие, из \ref{eq:AB=<C} и \ref{ex:hoelder-norm} следуют неравенства из \ref{lem:nash-guenter:guenter};
в частности, коэффициент перед $\|w\|_{C^{2,\alpha}}\cdot\|w\|_{C^{k,\alpha}}$ тот же при всех $k$.
\qeds

Обещанное касательное поле $\vec x=\vec x(w)$ можно определить как линейную комбинацию $\partial_if$ такую, что
\[\langle \partial_if,\vec x\rangle=-(\Delta-1)^{-1} A_i(w)\]
для каждого $i$; сравните со второй строкой в \ref{eq:tildeQ}.

\paragraph{Благодарности.}
Эта работа была частично поддержана Национальным научным фондом США, грант DMS-2005279.

Я благодарен
анонимному рецензенту,
Михаилу Громову,
Фёдору Назарову
и Дину Янгу за помощь.
